The exact response-type game is a strict treatment-domain extension of the unrestricted binary decision problem from two arms to every fixed \(K\geq3\); it is not a setwise extension of Hull's bounded-outcome model. When \(K=2\) and \(c=(1,-1)\), the orbit experiment retains the four distinct counts \(m_{(0,0)},m_{(0,1)},m_{(1,0)},m_{(1,1)}\) in \(P_m(x\mid r)\). Although
\[
 \tau_c(m)=\frac{m_{(1,0)}-m_{(0,1)}}n=\frac{p_+-p_-}{n},
\]
no aggregation by individual-effect class is asserted to identify the likelihood kernel. The exact orbit equality and the internal sign-group and Bayesian-information arguments give
\[
 \rho_n((1,-1))=G_n((1,-1)),
 \qquad C_0((1,-1))=1,
 \qquad d_n((1,-1))=\Theta(n^{-4/3}),
\]
with the explicit coarse ceiling \(d_n((1,-1))\leq43n^{-4/3}\) for \(n\geq8\), but no exact second-order constant or source-specific optimizer is delivered. Hull's theorem permits unrestricted designs and measurable estimators for two bounded-outcome arms. By contrast, for every fixed \(K\geq3\) and nonzero zero-sum \(c\), the present theorem assumes binary response schedules and proves equality of the unrestricted labeled-schedule risk and the explicit joint \(2^K\)-response-type orbit game, including arbitrary assignment dependence and possibly biased estimators. Thus the strict extension is on the number-of-arms axis within the binary model. This statement asserts neither a bounded-outcome \(K\geq3\) extension nor that the multi-arm game is scalar-nonrepresentable.